\section{多元函数微分学}
\subsection{多元函数的概念、极限与连续性}
\begin{mydef}{二元函数}{two-variable function}
    设\(D\)是平面上的一个点集，如果对每个点\(P(x, y) \in D\)，都有一个唯一的数\(z = f(x, y)\)，则称\(f(x, y)\)是\(D\)上的一个二元函数。
    \(D\)为该函数的定义域，所有函数值组成的集合\(f(D)=\{f(x,y):(x,y)\in D\}\)为值域。
\end{mydef}
\begin{conclusion}{二元函数的几何意义}{two-variable function geometry}
    空间点集\(\{(x,y,z):(x,y)\in D,\ z=f(x,y)\}\)为二元函数\(z=f(x,y)\)的图形。通常它是一张曲面。
    曲面\(z = f(x, y)\)与平面\(z = c\)的交线为\(z = f(x, y) = c\)，这是一个曲线。
    曲线\(z = f(x, y) = c\)在\(Oxy\)平面上的投影曲线:\(f(x, y) = C\)称为\(z = f(x, y)\)的\textbf{等高线}
\end{conclusion}
\begin{conclusion}{一元函数与二元函数的联系与区别}{one-variable function and two-variable function}
    \begin{enumerate}
        \item 一元函数是二元函数的特殊情形：让一自变量动，另一自变量固定，或让\((x, y)\)沿某曲线变动，二元函数就转换为一元函数
        \item 一元函数中，自变量\(x\)代表直线上的点，只有两个变动方向，而二元函数中，自变量\(x, y\)代表平面上的点，有无数个变动方向。
        \item 一元函数\(z = f(x)(a < x <b)\)也可看成二元函数，其定义域为\(
            a<x<b,\ -\infty<y<+\infty
        \)
    \end{enumerate}
\end{conclusion}
\begin{mydef}{二元函数极限}{two-variable function limit}
    设函数\(f(x, y)\)在开区域或闭区域\(D\)内有定义，\(M_0(x_0, y_0)\)是\(D\)的内点或边界点
    则\\
    \(
        \lim\limits_{\substack{x \to x_0 \\ y \to y_0}} f(x,y)= 
        \lim\limits_{(x, y) \to (x_0, y_0)}
            f(x, y) = A  \Leftrightarrow
        \forall \varepsilon > 0, \exists \delta > 0,
        \text{当}(x,y) \in D, \text{且}0 < \sqrt{(x - x_0)^2 + (y - y_0)^2} < \delta
        \text{时，有}\left| f(x, y) - A \right| < \varepsilon
    \)\\
    这里的极限过程是点\(x, y\)在\(D\)内沿任何路径以任何方式趋于点\((x_0, y_0)\)
\end{mydef}
\begin{conclusion}{二元函数极限与一元函数极限比较}{two-variable function limit and one-variable function limit}
    \begin{enumerate}
        \item 二元函数与一元函数有相同的极限运算法则与极限性质
        \\
        求二元函数极限的常用方法有
        \begin{enumerate}
            \item 用极限的运算法则
            \item 放大缩小法，变量变换法转为求简单的极限或一元函数的极限
        \end{enumerate}
        \item 二元函数\(z=f(x, y)\)与一元函数之极限存在性问题的区别
        \\
        与一元函数极限中"函数在某一点的极限当且仅当它在该点的左右极限存在且相等"相对应。\\
        在二元函数的极限中，极限 \(\lim\limits_{(x, y) \to (x_0, y_0)}
            f(x,y)=A\)的含义是：对定义域内任意趋于\((x_0,y_0)\)的点列（不取极限点本身），相应函数值都趋于\(A\)。只检查直线或少数几类路径只能用于否定极限，通常不能据此证明极限存在。\\
            \\
            若在定义域内沿两条不同路径\((\text{如},\text{直线}y - kx, \text{抛物线}y = x^2)\),极限\(
                \lim\limits_{(x, y) \to (x_0, y_0)} f(x, y)
            \)的值不相等或沿某一路径极限\(\lim\limits_{(x,y) \to (x_0, y_0)} f(x, y)\)不存在，则\(
                \lim\limits_{(x, y) \to (x_0, y_0)} f(x, y)
            \)不存在。这是证明二元函数极限不存在的一个有效方法
    \end{enumerate}
\end{conclusion}
\begin{mydef}{二元函数连续性}{two-variable function continuity}
    设\(z = f(x, y)\)定义在区域\(D\)上，\(P_0(x_0, y_0)\)是\(D\)的内点或边界点，若\(
        \lim\limits_{(x, y) \to (x_0, y_0)} f(x, y) = f(x_0, y_0)
    \)，则称\(f(x,y)\)在\(P_0(x_0,y_0)\)处连续。若对每个\((x_0,y_0)\in D\)上述条件都成立（边界点按定义域内趋近），则称\(f\)在\(D\)上连续。
\end{mydef}
\begin{conclusion}{二元连续函数的性质}{two-variable function continuity property}
    由自变量\(x\)的初等函数及自变量\(y\)的初等函数经过有限次四则运算和复合运算得到的二元函数
        \(z = f(x, y)\)称为二元初等函数，类似地可定义多元初等函数。\textbf{多元初等函数在其定义区域上是连续的} \\
        判断二元函数连续性有和一元函数相同的方法 \\
        二元连续函数具有以下相应性质
        \begin{enumerate}
            \item 设\(z = f(x,y )\)在\(M(x_0, y_0)\)连续，\(f(x_0, y_0) > 0(< 0)\),则
            \(\exists \delta > 0, \text{当}0 < \sqrt{(x - x_0)^2 + (y - y_0)^2} < \delta
            \text{且} (x, y) \in D \)时，\(f(x, y) > 0(< 0)\)
            \item 最大值最小值定理 \\
            设\(z = f(x, y)\)在有界闭区域\(D\)上连续，则它在\(D\)上一定有最大值和最小值，
            即\(\exists M_1,M_2 \in D,\)使得\(\forall M \in D\),有
            \[
                f(M_2) \leq f(M) \leq f(M_1)
            \]
            其中\(f(M_2)\)为最小值，\(f(M_1)\)为最大值
            \item 中间值定理 \\
            设\(z = f(x, y)\)在连通的区域\(D\)上连续，\(\forall M_1, M_2 \in D\) ，
            \(
                f(M_1) < f(M_2)
            \)，则对任意实数\(\mu\)满足\(f(M_1)<\mu<f(M_2)\)，
            则存在一点\(P(x, y)\)，使得\(f(P) = \mu\)
        \end{enumerate}
\end{conclusion}
\subsection{多元函数的偏导数与全微分}
\begin{mydef}{偏导数}{partial derivative}
    设二元函数\(z=f(x,y)\)在点\((x_0,y_0)\)的某邻域内有定义。固定其中一个变量后，若相应的一元函数导数存在，则称其为\(f\)在\((x_0,y_0)\)处关于另一个变量的偏导数，记为\\
    \(
        \frac{\partial f(x_0, y_0)}{\partial x}, \frac{\partial z}{\partial x} |_{(x_0, y_0)},
        \frac{\partial f}{\partial x} |_{(x_0, y_0)},z^{'}_{x} |_{(x_0, y_0)},\text{或}f^{'}_{x} |_{(x_0, y_0)}
    \)\\
    \(
        \left(
            \frac{\partial f(x_0, y_0)}{\partial y}, \frac{\partial z}{\partial y} |_{(x_0, y_0)},
        \frac{\partial f}{\partial y} |_{(x_0, y_0)},z^{'}_{y} |_{(x_0, y_0)},\text{或}f^{'}_{y} |_{(x_0, y_0)} 
        \right)\)
    按定义有
    \[
        \begin{aligned}
            \frac{\partial f}{\partial x}(x_0,y_0)=\lim\limits_{\Delta x\to0}\frac{f(x_0+\Delta x,y_0)-f(x_0,y_0)}{\Delta x},\\
            \frac{\partial f}{\partial y}(x_0,y_0)=\lim\limits_{\Delta y\to0}\frac{f(x_0,y_0+\Delta y)-f(x_0,y_0)}{\Delta y}.
        \end{aligned}
    \]
    若\(z = f(x, y)\)在区域\(D\)的每一点都有偏导数，一般地说，它们仍是\(x, y\)的函数，称为对\(f(x, y)\)的\textbf{偏导函数}，简称\textbf{偏导数},
    记为\(
        \frac{\partial z}{\partial x}, \frac{\partial f}{\partial x},f^{'}_{x}(x, y);
        \frac{\partial z}{\partial y}, \frac{\partial f}{\partial y},f^{'}_{y}(x, y);
    \)
    偏导数本质上是一元函数的导数，在计算时可以代入已知值，减少计算量
\end{mydef}


\begin{conclusion}{偏导数的几何意义}{partial derivative geometric meaning}
    \(f_x(x_0,y_0)\)是曲面\(z=f(x,y)\)与平面\(y=y_0\)的交线在\(M_0(x_0,y_0,f(x_0,y_0))\)处沿\(x\)方向的切线斜率；
    \(f_y(x_0,y_0)\)是曲面与平面\(x=x_0\)的交线在该点沿\(y\)方向的切线斜率。
\end{conclusion}
\begin{conclusion}{偏导数的计算}{partial derivative calculation}
    求偏导数，归结为求一元函数的导数
    求\(f(x, y) = \begin{cases}
            g(x,y), (x, y) \neq (x_0, y_0) \\
            A, (x, y) = (x_0, y_0)
        \end{cases}\)
        在\((x_0, y_0)\)处偏导数的方法
    \begin{enumerate}
    \item  按定义
    \[
    \begin{aligned}
        \frac{\partial f(x_0, y_0)}{\partial x} = 
    \lim\limits_{\varDelta x \to 0} \frac{f(x_0 + \varDelta x, y_0) - f(x_0, y_0)}{\varDelta x} \\
    = \lim\limits_{\varDelta x \to 0} \frac{g(x_0 + \varDelta x, y_0) - A}{\varDelta x}
    \end{aligned}
    \]
    类似可求偏导数\(\frac{\partial f(x_0, y_0)}{\partial y}\)
    \item 利用一元函数的导数极限定理
    若一元函数\(x\mapsto f(x,y_0)\)在\(x_0\)处连续、在\(x_0\)的去心邻域内可导，且
    \[
      \lim_{x\to x_0}f_x(x,y_0)=B,
    \]
    则由拉格朗日中值定理可得\(f_x(x_0,y_0)=B\)。关于\(y\)的偏导数同理。注意应取邻域内偏导数的极限，而不是函数值\(g(x,y_0)\)的极限。
    \end{enumerate}
\end{conclusion}
\begin{mydef}{可微性与全微分}{differentiability and total differential}
    设二元函数\(z = f(x, y)\)在点\((x_0, y_0)\)的某邻域内有定义，如果该函数在点\(M_0(x_0, y_0)\)处的全增量\(
    \varDelta z = f(x_0 + \varDelta x, y_0 + \varDelta y) - f(x_0, y_0)\),可表示为\[
        \varDelta z = A \varDelta x + B \varDelta y + o(\sqrt{\varDelta x^2 + \varDelta y^2})
    \]
    可以将\(\sqrt{\varDelta x^2 + \varDelta y^2}\)记为\(\rho(\rho \to 0)\)，其中\(A,B\)不依赖于\(
        \varDelta x \varDelta y,
    \)而仅与点\((x_0, y_0)\)有关，则称\(z = f(x, y)\)在点\((x_0, y_0)\)处可微,而\(A\varDelta x + B \varDelta y\)称为\(z = f(x, y)\)在点\((x_0, y_0)\)处的全微分，记为\[
        \mathrm{d}z |_{(x_0, y_0)} = A \varDelta x + B \varDelta y
    \]
\end{mydef}
\begin{conclusion}{$z = f(x, y)$在点$(x_0, y_0)$可微分的必要与充分条件}{all derivative requirement}
    \begin{enumerate}
        \item 可微的必要条件 \\
            若\(f(x, y)\)在点\((x_0, y_0)\)可微，则\(z = f(x, y)\)在点\((x_0, y_0)\)处连续，两个偏导数\(
                f'_{x}(x_0, y_0), f'_{y}(x_0, y_0)
            \)都存在,且\(A = f'_{x}(x_0, y_0), B = f'_{y}(x_0, y_0)\)即有
            \[
                \mathrm{d}z |_{(x_0, y_0)} = f'_{x}(x_0, y_0) \mathrm{d} x + f'_{y}(x_0, y_0) \mathrm{d} y
            \]
            其中\(\mathrm{d}x = \varDelta x, \mathrm{d} y = \varDelta y\)
        \item 可微的充分条件 \\
        若\(z = f(x, y)\)在点\((x_0, y_0)\)的某邻域内存在偏导数\(f'_{x}(x, y), f'_{y}(x, y)\)，且这两个偏导数在点\((x_0, y_0)\)处连续，则\(z = f(x, y)\)在点\((x_0, y_0)\)处可微
    \end{enumerate}
\end{conclusion}

\begin{conclusion}{偏导数的连续性，函数可微性，可偏导性与函数连续性的关系}{inter relationship}
    若\(z = f(x, y)\)在点\((x_0, y_0)\)处存在\(f'_{x}(x_0, y_0), f'_{y}(x_0, y_0)\)则称\(f(x, y)\)在点\((x_0, y_0)\)处可偏导。
    两个偏导数连续，可微，可偏导于二元函数连续的关系为 \\
    %一个树形的关系
    \begin{tikzpicture}
        \node[rectangle, draw] (continuous) {\(f'_{x}(x, y),f'_{y}(x, y)\)在点\((x_0,y_0)\)处连续};
        \node[rectangle, draw, below = of continuous] (partial) {
            \(z = f(x, y)\)在点\((x_0,y_0)\)处可微且\(dz = Adx + Bdy\);
        };
        \node[rectangle, draw, below = of partial] (continuousf) {
            \(f(x,y)\)在点\((x_0,y_0)\)处连续
        };
        \node[rectangle, draw, below = of continuousf] (derivative) {
            \(z = f(x, y)\)在点\((x_0,y_0)\)处可偏导且\(f'_{x}(x_0, y_0) = A, f'_{y}(x_0, y_0) = B\)
        };
        %连线，线使用\Leftarrow类似的箭头
        \draw[double, -Implies, double distance=1pt] (continuous) -- (partial);
        %线在左边
        \draw[double, -Implies, double distance=1pt] ([xshift=-82pt]partial.east) -- ([xshift=-10pt]continuousf.east);
        \draw[double, -Implies, double distance=1pt] (partial.west) -- ([xshift=42pt] derivative.west);
        \draw[double, -Implies, double distance=1pt] (continuousf.west) to node[sloped, allow upside down] {//} ([xshift=110pt]derivative.west);
        \draw[double, -Implies, double distance=1pt] ([xshift=-110pt]derivative.east) to node[sloped, allow upside down] {//} (continuousf.east);
        \draw[double, -Implies, double distance=1pt] (continuousf.west) to node[sloped, allow upside down] {//} ([xshift=80pt] partial.west);
        \draw[double, -Implies, double distance=1pt] (derivative.east) to node[sloped, allow upside down] {//} (partial.east);
        \draw[double, -Implies, double distance=1pt] (partial) to node[sloped, allow upside down] {//} ([xshift=10pt]continuous);
    \end{tikzpicture}\\
    在一元函数中，可导和可微是等价的，且都强于连续。\\
    对于多元函数，情况变得复杂：可微最强，它同时蕴含连续和所有方向的可偏导；但即便所有方向都可偏导，也推不出可微，甚至推不出连续。
\end{conclusion}

\begin{mydef}{高阶偏导数、混合偏导数}{higher order partial derivative and mixed partial derivative}
    若函数\(z = f(x, y)\)的一阶偏导数\(\frac{\partial z}{\partial x} = f'_{x}, \frac{\partial z}{\partial y} = f'_{y}\)关于\(x,y\)的偏导数仍然存在。
    则称一阶偏导数的偏导数是\(z = f(x, y)\)的二阶偏导。根据求导顺序的不同有以下四个二阶偏导数
    \[
        \begin{aligned}
            \frac{\partial^2 z}{\partial x^2} = \frac{\partial}{\partial x}\left(\frac{\partial z}{\partial x}\right)  = f''_{xx}(x, y);\quad
            \frac{\partial^2 z}{\partial x \partial y} = \frac{\partial}{\partial y}\left(\frac{\partial z}{\partial x}\right)  = f''_{xy}(x, y) \\
             \frac{\partial^2 z}{\partial y^2} = \frac{\partial}{\partial y}\left(\frac{\partial z}{\partial y}\right)  = f''_{yy}(x, y);\quad
            \frac{\partial^2 z}{\partial y \partial x} = \frac{\partial}{\partial x}\left(\frac{\partial z}{\partial y}\right)  = f''_{yx}(x, y) \\
        \end{aligned}
    \]
    \(f_{xy}\)与\(f_{yx}\)称为混合偏导数；二者在没有附加条件时未必相等。
    同理可得更高阶的偏导数。
\end{mydef}

\begin{conclusion}{高阶偏导数、混合偏导数次序无关}{order irrelevant}
    若二元函数\(z = f(x, y)\)的两个二阶混合偏导数\(
         \frac{\partial^2 z}{\partial y \partial x},
         \frac{\partial^2 z}{\partial x \partial y}
    \)都在区域\(D\)上连续，则在区域\(D\)内\(
        \frac{\partial^2 z}{\partial y \partial x} =
        \frac{\partial^2 z}{\partial x \partial y}
    \),即二阶混合偏导数与求偏导的先后次序无关。
\end{conclusion}


\begin{conclusion}{多元函数为常数的条件}{constant function}
    \begin{enumerate}
        \item 设\(D\)是连通区域，\(f(x,y)\)在\(D\)上可微且满足\(\frac{\partial f}{\partial x}=0,\frac{\partial f}{\partial y}=0(
            \forall (x, y) \in D
        )\),则在区域\(D\)内，\(f(x, y)\)为常数。\\又在区域\(D\)上\(df(x,y) = 0 \Leftrightarrow\)在区域\(D\)上\(
            \frac{\partial f}{\partial x} = \frac{\partial f}{\partial y} = 0,
        \)因此也有\\
        设\(f(x,y)\)在区域\(D\)上满足\(\mathrm{d}f(x, y) = 0\),则\(f(x, y)\)在区域\(D\)为常数
        \item 设\(z = f(x, y)\)定义在全平面上，若\(\frac{\partial f}{\partial x} = 0\),则\(
            f(x, y) = \varphi(y);
        \)同理，若\(\frac{\partial f}{\partial y} = 0\),则\(f(x, y) = \psi(x)\)。\\
        其中\(\psi(x)\)是\(x\)的函数，\(\varphi(y)\)是\(y\)的函数。二者都是一元函数。
    \end{enumerate}
\end{conclusion}

\subsection{多元函数的微分法与复合函数求导法则}
\begin{formula}{全微分四则运算法则}{total differential formula}
    若\(u = u(x, y), v = v(x, y)\)在点\((x, y)\)处可微，则
    \[
        \begin{aligned}
            &\mathrm{d}(u \pm v) = \mathrm{d}u \pm \mathrm{d}v;\quad 
            \mathrm{d}(cu) = c\,\mathrm{d}u\quad(c\text{为常数})\\[1em]
            &\mathrm{d}(u \cdot v) = u \mathrm{d}v + v \mathrm{d}u; \quad
            \mathrm{d}\left(\frac{u}{v}\right) = \frac{v\mathrm{d}u  - u \mathrm{d}v}{v^2};
        \end{aligned}
    \]
\end{formula}
\begin{conclusion}{多元复合函数的微分法则}{multi-variable function derivative rule}
    \begin{enumerate}
    \item 多元函数与一元函数复合 \\
    设\(x = x(t), y = y(t), z = z(t)\)在\(t\)可导, \(u = f(x, y, z)\)在对应点\((x, y, z) = 
    (x(t), y(t), z(t))\)可微，则\(u = f(x(t), y(t), z(t))\)在点\(t\)可导，且
    \[
        \frac{\mathrm{d}u}{\mathrm{d}t} = \frac{\partial u}{\partial x}\frac{\mathrm{d}x}{\mathrm{d}t} + 
        \frac{\partial u}{\partial y}\frac{\mathrm{d}y}{\mathrm{d}t} + 
        \frac{\partial u}{\partial z}\frac{\mathrm{d}z}{\mathrm{d}t};
    \]
    在这里, \(\frac{\mathrm{d}u}{\mathrm{d}t}\)也被称为全导数
    \item 多元函数与多元函数的复合 \\
    \begin{enumerate}
        \item 链式法则 \\
        设\(u=u(x,y),v=v(x,y)\)在点\((x,y)\)可微，\(z=f(u,v)\)在对应点可微，则复合函数\(z=f(u(x,y),v(x,y))\)满足
        \[
            \frac{\partial z}{\partial x} = \frac{\partial f}{\partial u}\frac{\partial u}{\partial x} + 
            \frac{\partial f}{\partial v}\frac{\partial v}{\partial x};\quad 
            \frac{\partial z}{\partial y} = \frac{\partial f}{\partial u}\frac{\partial u}{\partial y} + 
            \frac{\partial f}{\partial v}\frac{\partial v}{\partial y};
        \]
        类似地，设\(u=u(x,y),v=v(x,y),w=w(x,y)\)在点\((x,y)\)可微，且\(z=f(u,v,w)\)在对应点可微，则
        \[
            \frac{\partial z}{\partial x} = \frac{\partial f}{\partial u}\frac{\partial u}{\partial x} + 
            \frac{\partial f}{\partial v}\frac{\partial v}{\partial x} + 
            \frac{\partial f}{\partial w}\frac{\partial w}{\partial x};\quad 
            \frac{\partial z}{\partial y} = \frac{\partial f}{\partial u}\frac{\partial u}{\partial y} + 
            \frac{\partial f}{\partial v}\frac{\partial v}{\partial y} + 
            \frac{\partial f}{\partial w}\frac{\partial w}{\partial y};
        \]
        \\
        为了方便，我们常用\(f'_{1}\)表示对第一个变量求偏导数,
        类似地有\(
           f'_{2} = \frac{\partial f}{\partial v}\)表示对第二个变量求偏导数，\(f'_{3} = \frac{\partial f}{\partial w}\)表示对第三个变量求偏导数。
           \\
           同理，\(f'_{11}\)表示对第一个变量求二阶偏导数，\(f'_{22}\)表示对第二个变量求二阶偏导数，\(f'_{33}\)表示对第三个变量求二阶偏导数。\\
           \\
           \(f'_(12)\)表示对第一个变量求偏导数，再对第二个变量求偏导数
        \item 一阶全微分不变性 \\
        设函数\(z = f(u, v), u = \phi(x, y), v = \psi(x, y)\)在点\((x, y)\)处可微，则
        可得函数\(z= f(u, v) = f(\phi(x, y), \psi(x, y))\)的全微分为
        \[
            \begin{aligned}
            \mathrm{d}z &= \frac{\partial z}{\partial x} \mathrm{d}x + \frac{\partial z}{\partial y} \mathrm{d}y \\
            & = \frac{\partial f}{\partial u} \left(
                \frac{\partial u}{\partial x} \mathrm{d}x + \frac{\partial u}{\partial y} \mathrm{d}y
            \right)
            + 
            \frac{\partial f}{\partial v} \left(
                \frac{\partial v}{\partial x} \mathrm{d}x + \frac{\partial v}{\partial y} \mathrm{d}y
                \right)\\
            &=  
            \frac{\partial f}{\partial u} \mathrm{d}u + 
            \frac{\partial f}{\partial v} \mathrm{d}v;
            \end{aligned}
        \]
        这表明当\(f, u, v\)都是可微函数时，尽管\(u, v\)不是自变量,但函数\(
            z = f(u, v) 
        \)的一阶全微分也具有与\(u, v\)是自变量时的一阶全微分的相同形式，这种性质称为\textbf{一阶全微分的不变性}。 \\
        在计算中，可以先求出\(u, v\)的全微分，再代入\(z = f(u, v)\)对\(u,v\)中计算全微分。
        \\
    \end{enumerate}
    \end{enumerate}
\end{conclusion}
\begin{conclusion}{复合函数求导法则的注意事项}{compound function derivative rule notice}
    \begin{enumerate}
        \item 弄清复合函数的结构，分清中间变量与自变量
        \item 偏导数公式的构成，复合函数对指定自变量的偏导数，求中间变量时不要漏项
        \item 在复合函数求导公式中，函数对中间变量的偏导数仍是中间变量的函数。如设\(
            z = f(u, v), u = \phi(x, y), v = \psi(x, y)\)均具有计算所需的连续偏导数，则
            \(
                \frac{\partial z}{\partial x} = \frac{\partial f}{\partial u}\frac{\partial u}{\partial x} +
                \frac{\partial f}{\partial v}\frac{\partial v}{\partial x}.
            \)
            这里\(\frac{\partial f}{\partial u}\)和\(\frac{\partial f}{\partial v}\)复合后也是\(x,y\)的函数；对\(\frac{\partial z}{\partial x}\)继续求导时，仍要使用链式法则：
            \[
            \frac{\partial}{\partial x}\left(\frac{\partial f}{\partial u}\right)
                 = \frac{\partial^2 f}{\partial u^2}\frac{\partial \phi}{\partial x}
                 + \frac{\partial^2 f}{\partial v\partial u}\frac{\partial \psi}{\partial x}
            \]
            \newline
            \[
                \frac{\partial}{\partial x}\left(\frac{\partial f}{\partial v}\right)
                = \frac{\partial^2 f}{\partial u\partial v}\frac{\partial \phi}{\partial x}
                + \frac{\partial^2 f}{\partial v^2}\frac{\partial \psi}{\partial x}
            \]
            \newline
            \[
                \frac{\partial^2 z}{\partial x^2} = \frac{\partial}{\partial x}\left(
                    \frac{\partial f}{\partial u}
                    \right)\frac{\partial \phi}{\partial x} +
                \frac{\partial f}{\partial u}\frac{\partial^2 \phi}{\partial x^2}
                + \frac{\partial}{\partial x}\left(
                    \frac{\partial f}{\partial v}
                    \right)\frac{\partial \psi}{\partial x} +
                    \frac{\partial f}{\partial v}\frac{\partial^2 \psi}{\partial x^2}
            \]\\
            将\(
            \frac{\partial}{\partial x}\left(\frac{\partial f}{\partial v}\right)\)
            与
            \(
            \frac{\partial}{\partial x}\left(\frac{\partial f}{\partial u}\right)
            \)代入上式，可得最后结果
    \end{enumerate}
\end{conclusion}

\begin{conclusion}{由一个方程式确定的一元隐函数求导法(隐函数存在定理)}{implicit-function-one-variable}
    如果二元方程\(F(x, y) = 0\)满足如下三个条件
    \begin{enumerate}
        \item 函数\(F(x, y)\)在点\((x_0, y_0)\)某邻域有连续的偏导数
        \item \(F(x_0, y_0) = 0\)
        \item \(F'_{y}(x_0, y_0) \neq 0\)，则在点\((x_0, y_0)\)附近，方程\(F(x, y) = 0\)有唯一解\(y=y(x)\)。且
        \[
            y' = \frac{\mathrm{d}y}{\mathrm{d}x} = - \frac{F'_x}{F'_y} 
        \]
    \end{enumerate}
\end{conclusion}
\begin{conclusion}{由一个方程式确定的二元隐函数求导法}{implicit-function-two-variable}
    二元隐函数的隐函数存在定理为:
    如果三元方程\(F(x, y, z) = 0\)满足如下三个条件
    \begin{enumerate}
        \item 函数\(F(x, y, z)\)在点\((x_0, y_0, z_0)\)某邻域有连续的偏导数
        \item \(F(x_0, y_0, z_0) = 0\)
        \item \(F'_{z}(x_0, y_0, z_0) \neq 0\)，则在点\((x_0, y_0, z_0)\)附近，方程\(F(x, y, z) = 0\)有唯一解，满足
        \(z_0 = z(x_0, y_0)\)并有连续的偏导数，且有
        \(
            \frac{\partial z}{\partial x} = - \frac{F'_x}{F'_z};
            \frac{\partial z}{\partial y} = - \frac{F'_y}{F'_z},
        \)
    \end{enumerate}
    求二元隐函数的偏导数的方法有
    \begin{enumerate}
        \item 求由方程\(F(x, y, z) = 0\)确定的隐函数\(z = z(x, y)\)的一阶偏导数求导方法 \\
        \begin{enumerate}
            \item 法一 \\
            由\(
                F(x, y, z(x, y)) = 0
            \)
            分别对\(x,y\)求偏导数，利用复合函数求导法得
            \[
                F'_x + F'_z\frac{\partial z}{\partial x}=0,\qquad
                F'_y + F'_z\frac{\partial z}{\partial y}=0.
            \]
            分别解出\(\frac{\partial z}{\partial x}, \frac{\partial z}{\partial y}\)
            \item 法二 \\
            利用一阶全微分不变性，对方程\(F(x, y, z) = 0\)两边同时对全微分得到\(
                \mathrm{d}F(x, y, z) = 0,
                F'_x\mathrm{d}x + F'_y\mathrm{d}y + F'_z\mathrm{d}z = 0
            \)
            从而当\(F'_z \neq 0\)时有\(
                \mathrm{d}z = - \left(\frac{F'_x}{F'_z}\mathrm{d}x
                + \frac{F'_y}{F'_z}\mathrm{d}y\right)
                \)
            于是有\(
                \frac{\partial z}{\partial x} = - \frac{F'_x}{F'_z};
                \frac{\partial z}{\partial y} = - \frac{F'_y}{F'_z},
            \)
            \item 法三 \\
            公式法求出\(
                \frac{\partial z}{\partial x} = - \frac{F'_x}{F'_z};
                \frac{\partial z}{\partial y} = - \frac{F'_y}{F'_z},
            \)
        \end{enumerate}
        \item 求由方程\(F(x, y, z) = 0\)确定的隐函数\(z = z(x, y)\)二阶偏导数求导方法 \\
        若函数\(F(x, y, z)\)具有二阶偏导数，则隐函数\(z = z(x, y)\)的二阶偏导数存在,可求\(z=z(x, y)\)的二阶偏导数，
        此时\(
            F'_x, F'_y, F'_z\)仍然是\(x, y, z\)的函数,同时也是\(x, y\)的函数。
            此处,使用\(F'_1, F'_2, F'_3\)代替\(F'_x, F'_y, F'_z\)。
        \begin{enumerate}
            \item 法一 \\
            将\(F'_x + F'_z\frac{\partial z}{\partial x}\)两边同时对x求偏导数，可得
            \[
                \frac{\partial}{\partial x}(F'_1) + \frac{\partial}{\partial x}(F'_3)\frac{\partial z}{\partial x} +
                F'_3\frac{\partial^2 z}{\partial x^2} = 0
            \]
            得
            \[
                \frac{\partial^2 z}{\partial x^2} = - \left[
                    F''_{11} + 2F''_{13}\frac{\partial z}{\partial x} + F''_{33}\left(\frac{\partial z}{\partial x}\right)^2
                \right] / F'_3
            \]
            整理得
            \[
                \frac{\partial^2 z}{\partial x^2} = - \left[
                    F''_{11}{F'_{3}}^2 -
                    2F''_{13}F'_{1}F'_{3} +
                    F''_{33}{F'_{1}}^2
                    \right] / {F'_{3}}^3
            \]
            \item 法二 \\
            对\(
                \frac{\partial z}{\partial x} = - \frac{F'_x}{F'_z};
            \)继续求偏导数，可得
            \[
                \frac{\partial^2 z}{\partial x^2} = -\left[
                    \left(F''_{11} + F''_{13}\frac{\partial z}{\partial x}\right)F'_{3} -F'_1 \left(F''_{31} + F''_{33}\frac{\partial z}{\partial x}
                    \right)\right] / {(F'_{3})}^{2}
            \]
            将\(
            \frac{\partial z}{\partial x} = - \frac{F'_1}{F'_3};
            \)代入，也可以得到法一的结果
        \end{enumerate}
    \end{enumerate}
\end{conclusion}

\begin{conclusion}{由方程组确定的一元隐函数求导法}{implicit-system-one-variable}
    \(
    \begin{cases}
        F(x, u, v) = 0 \\
        G(x, u, v) = 0 \
    \end{cases}
    \)其中\(F, G\)有连续偏导数，若在区间\(I\)上存在函数组\(
        u = u(x), v = v(x)
    \)满足方程组
    应用复合函数求导法则，可得
    \[
        \begin{cases}
            F'_{1} + F'_{2}\frac{\partial u}{\partial x} + F'_{3}\frac{\partial v}{\partial x} = 0 \\

            G'_{1} + G'_{2}\frac{\partial u}{\partial x} + G'_{3}\frac{\partial v}{\partial x} = 0 \
        \end{cases}
    \]
    将其看成是以\(\frac{\mathrm{d}u}{\mathrm{d}x}, \frac{\mathrm{d}v}{\mathrm{d}x}\)为未知数的二元一次方程组，利用克莱姆法则可解
\end{conclusion}
\begin{conclusion}{由方程组确定的二元隐函数求导法}{implicit-system-two-variable}
    设有\(4\)个变量和\(2\)个方程构成的方程组\(
        \begin{cases}
            F(x, y, u, v) = 0 \\
            G(x, y, u, v) = 0 
        \end{cases}
    \)，设\(F,G\)有连续偏导数，且
    \[
        \frac{\partial(F,G)}{\partial(u,v)}
        =\begin{vmatrix}F_u&F_v\\G_u&G_v\end{vmatrix}\ne0.
    \]
    则该方程组在相应点附近确定可微隐函数\(u=u(x,y),v=v(x,y)\)。分别对\(x,y\)求偏导，得
    \[
        \begin{cases}
            F_x + F_u\frac{\partial u}{\partial x} + F_v\frac{\partial v}{\partial x} = 0 \\
            G_x + G_u\frac{\partial u}{\partial x} + G_v\frac{\partial v}{\partial x} = 0
        \end{cases}
        \text{或}
        \begin{cases}
            F_y + F_u\frac{\partial u}{\partial y} + F_v\frac{\partial v}{\partial y} = 0 \\
            G_y + G_u\frac{\partial u}{\partial y} + G_v\frac{\partial v}{\partial y} = 0
        \end{cases}
    \]
    把它们分别看成关于\(u_x,v_x\)或\(u_y,v_y\)的二元一次方程组，由克莱姆法则
    \[
    \begin{aligned}
        u_x&=-\frac{\partial(F,G)/\partial(x,v)}{\partial(F,G)/\partial(u,v)},&
        v_x&= \frac{\partial(F,G)/\partial(x,u)}{\partial(F,G)/\partial(u,v)},\\
        u_y&=-\frac{\partial(F,G)/\partial(y,v)}{\partial(F,G)/\partial(u,v)},&
        v_y&= \frac{\partial(F,G)/\partial(y,u)}{\partial(F,G)/\partial(u,v)}.
    \end{aligned}
    \]
\end{conclusion}
\begin{conclusion}{复合函数求导法则的其他应用}{compound function derivative rule other application}
    做变量替换时，求函数在新变量下得偏导数，有时通过变量替换可将某些偏微分方程化简
\end{conclusion}
\begin{conclusion}{极坐标变换下得拉普拉斯方程}{laplace equation under polar coordinate transformation}
    在极坐标变换下，拉普拉斯方程
    \[
        \frac{\partial^2 u}{\partial x^2} + \frac{\partial^2 u}{\partial y^2} = 0
    \]
    为
    \[
        \frac{\partial^2 u}{\partial r^2} + \frac{1}{r}\cdot\frac{\partial u}{\partial r} + \frac{1}{r^2}\cdot\frac{\partial^2 u}{\partial \theta^2} = 0
    \]
    其中\(u\)是二元函数。
\end{conclusion}
\begin{conclusion}{多元函数问题转化为一元函数问题}{multi-variable function problem to one-variable function}
    在求二元函数\(z = f(x,y)\)在\((x_0,y_0)\)处的二阶泰勒公式时，把二元函数的增量\(
        f(x_0 + \Delta x, y_0 + \Delta y) - f(x_0, y_0)
    \)看成一元函数\(\Phi(t) = f(x_0 + t\Delta x, y_0+t\Delta y)\)从\(t=0\)到\(t=1\)的增量\(
        \Phi(1) - \Phi(0)
    \)\\
    由一元函数的二阶泰勒公式，存在\(\theta\in(0,1)\)使
    \[
        \Phi(1) = \Phi(0) + \Phi'(0) + \frac{1}{2}\Phi''(0)
        + \frac{1}{3!}\Phi'''(\theta).
    \]
\end{conclusion}

\begin{formula}{二元函数的泰勒公式}{two-variable function taylor formula}
    设二元函数\(z = f(x, y)\)在点\(M_0(x_0, y_0)\)处的某邻域\(U(M_0)\)内具有三阶连续偏导数，\(
        (x_0 + \varDelta x, y_0 + \varDelta y) \in  U(M_0)
    \),则有二阶泰勒公式为
    \[
        \begin{aligned}
                   & f(x_0 + \varDelta x, y_0 + \varDelta y)  \\
                   &= f(x_0, y_0) + \left(\varDelta x\frac{\partial}{\partial x} + \varDelta y\frac{\partial}{\partial y}\right)f(x_0, y_0) \\
                    &+ \frac{1}{2!}\left(
                        \varDelta x\frac{\partial}{\partial x} + \varDelta y\frac{\partial}{\partial y}
                        \right)^2 f(x_0, y_0) + R_2 \\
                    &= f(x_0, y_0) + \frac{\partial f(x_0, y_0)}{\partial x} \varDelta x + \frac{\partial f(x_0, y_0)}{\partial y} \varDelta y \\
                    &+ \frac{1}{2}\left[
                        \frac{\partial^2 f(x_0, y_0)}{\partial x^2} \varDelta x^2 
                        + 2\frac{\partial^2 f(x_0, y_0)}{\partial x \partial y} \varDelta x \varDelta y
                        + \frac{\partial^2 f(x_0, y_0)}{\partial y^2} \varDelta y^2
                        \right] + R_2
        \end{aligned}
    \]
    这里\(R_2 = \frac{1}{3!}\left(
            \varDelta x\frac{\partial}{\partial x} + \varDelta y\frac{\partial}{\partial y}
        \right)^3 f(x_0 + \theta \varDelta x, y_0 + \theta \varDelta y) \ (0 < \theta < 1)
\)是拉格朗日余项。
若\(R_2 = o(\rho^2)\ (\rho = \sqrt{\varDelta x^2 + \varDelta y^2} \to 0)\),称为皮亚诺余项。\\
一般地，记\(D_{\Delta}=\varDelta x\,\partial_x+\varDelta y\,\partial_y\)。若\(f\)具有\(n+1\)阶连续偏导数，则存在\(\theta\in(0,1)\)使
\[
f(x_0+\varDelta x,y_0+\varDelta y)
=\sum_{m=0}^{n}\frac{1}{m!}D_{\Delta}^{m}f(x_0,y_0)
+\frac{1}{(n+1)!}D_{\Delta}^{n+1}
f(x_0+\theta\varDelta x,y_0+\theta\varDelta y),
\]
其中
\[
D_{\Delta}^{m}f
=\sum_{k=0}^{m}\binom{m}{k}
\frac{\partial^{m}f}{\partial x^k\partial y^{m-k}}
\varDelta x^k\varDelta y^{m-k}.
\]
\end{formula}
\subsection{多元函数的极值与最值问题}
\begin{mydef}{多元函数极值与驻点}{multi-variable function extremum and stationary point}
    若\(\exists M_0 (x_0, y_0)\)的某邻域\(U(M_0, \delta)\)使得
    \[
        f(x, y) \leq f(x_0, y_0) \qquad  (f(x, y) \geq f(x_0, y_0)), \forall (x, y) \in U(M_0, \delta)
    \]
    则称\(
    f(x, y)\)在\(M_0\)有极大值(极小值)\(f(x_0, y_0)\)。极大值与极小值点统称为极值
    \(M_0\)称为极值点。\\
    凡是能使\(f'_{x}(x, y) = 0, f'_{y}(x, y) = 0\)同时成立的点\(x, y\)称为\(
    z = f(x, y)\)的驻点
\end{mydef}
\begin{conclusion}{多元函数取得极值的必要条件与充分条件}{multi-variable function extremum necessary and sufficient condition}
    \begin{enumerate}
        \item 必要条件\\
        设\(f(x, y)\)在点\((x_0, y_0)\)处具有偏导数，且在点\((x_0, y_0)\)处取得极值，则
        它在该点的偏导数必然为零，即\(
            f'_{x}(x_0, y_0) = 0, f'_{y}(x_0, y_0) = 0
        \) \\
        \boxed{\text{注}} 具有偏导数的极值点必然是驻点，但驻点不一定是极值点。\\
        \item 充分条件 \\
        设\(z = f(x, y)\)在点\((x_0, y_0)\)的某邻域内连续且有一阶及二阶连续偏导数，又
        \(
        f'_{x}(x_0, y_0) = 0, f'_{y}(x_0, y_0) = 0
        \),令\(
            f''_{xx}(x_0,y_0)=A,\qquad f''_{xy}(x_0,y_0)=B,\qquad f''_{yy}(x_0,y_0)=C
        \),则
        \begin{enumerate}
            \item 若\(AC - B^2 > 0\),则\((x_0, y_0)\)是\(f(x, y)\)的极值点，且
            当\(A > 0\)时取得极小值，当\(A < 0\)时取得极大值。
            \item 若\(AC - B^2 < 0\),则\((x_0, y_0)\)不是\(f(x, y)\)的极值点
            \item 若\(AC - B^2 = 0\),\(f(x, y)\)在\((x_0, y_0)\)处可能取极值，也可能不取极值，还需
            另作讨论(一般转到极值定义法)
        \end{enumerate}
    \end{enumerate}
\end{conclusion}
\begin{conclusion}{求二元函数极值点的步骤}{two-variable function extremum point}
    若\(z = f(x, y)\)有连续的二阶偏导数，求极值点的步骤如下：
    \begin{enumerate}
        \item 求驻点，即解方程组\(
            \begin{cases}
                f'_{x}(x, y) = 0 \\
                f'_{y}(x, y) = 0
            \end{cases}
        \)
        \item 对每个驻点求出二阶导函数的值\(
            f''_{xx}(x_0, y_0) = A, 
            f''_{xy}(x_0, y_0) = B,
            f''_{yy}(x_0, y_0) = C
        \)
        \item 判断极值点，即根据\(AC - B^2\)的值来判断是否为极值点，若有，
        根据\(A\)来判断是极大值还是极小值。
    \end{enumerate}
\end{conclusion}
\begin{conclusion}{条件极值点的必要条件}{condition extremum point necessary condition}
    设\(f(x, y), \varphi(x, y)\)在点\((x_0, y_0)\)的某邻域内具有一阶连续偏导数，且
    \(\nabla\varphi(x_0,y_0)\neq\mathbf0\), 若\(
        P_0(x_0, y_0)\)是二元函数\(z = f(x, y)\)在条件\(\varphi(x, y) = 0\)下的极值点，则
    \[
        \begin{cases}
            \begin{vmatrix}
                f'_{x}(P_0) & f'_{y}(P_0) \\
                \varphi'_{x}(P_0) & \varphi'_{y}(P_0)
            \end{vmatrix} = 0 \\
            \varphi(P_0) = 0
        \end{cases}
        \Leftrightarrow \quad
        \begin{cases}
            f'_{x}(P_0) + \lambda \varphi'_{x}(P_0) = 0 \\
            f'_{y}(P_0) + \lambda \varphi'_{y}(P_0) = 0 \\
            \varphi(P_0) = 0
        \end{cases}
    \]
\end{conclusion}

\begin{conclusion}{极值最值问题的提法}{extremum and extremum point}
    \begin{enumerate}
        \item 简单极值(最值)问题\\
        求二元函数\(z = f(x, y)\)在\(
        Oxy\)平面上某区域\(D\)上的最大值或最小值 \\
        求三元函数\(u = f(x, y, z)\)在\(
        Oxyz\)空间中某区域\(\Omega\)上的最大值或最小值
        \item 条件极值(最值)问题\\
        求二元函数\(z = f(x, y)\)在条件\(\varphi(x, y) = 0\)下的最大值或最小值 \\
        求三元函数\(u = f(x, y, z)\)在条件\(\varphi(x, y, z) = 0\)或条件\(
            \phi(x, y, z) = 0, \psi(x, y, z) = 0
        \)下的最大值或最小值问题。
    \end{enumerate}
\end{conclusion}

\begin{conclusion}{求二元函数与三元函数的简单极值问题}{simple extremum problem}
    \begin{enumerate}
        \item 设\(z = f(x, y)\)在平面有界闭区域\(D\)上连续，且\(f(x, y)\)在\(D\)上存在最大值\(M\)与最小值\(m\)，它们
        或在\(D\)内满足\(
        \frac{\partial f}{\partial x} = 0, \frac{\partial f}{\partial y} = 0
        \)的点(即\(f(x, y)\)的驻点)取得，或在\(D\)的边界上取得，或在偏导数不存在的点上取得。
        或在\(D\)的边界上达到，因此求连续函数\(f(x, y)\)在有界闭区域\(D\)上的最值的步骤为
        \begin{enumerate}
            \item 求出\(f(x, y)\)在\(D\)内可能取得极值点(驻点和偏导数不存在的点)的函数值
            \item 求出\(f(x, y)\)在\(D\)的边界上的最大值与最小值
            \item 比较上述三个值，其中最大者为\(f(x, y)\)在\(D\)上的最大值，最小者 为\(f(x, y)\)在\(D\)上的最小值。
        \end{enumerate}
        \item 设\(z = f(x, y)\)在开区域\(D\)内可偏导,且根据实际问题可知它在\(D\)内有最大值或最小值，只需在满足\(
            \frac{\partial f}{\partial x} = 0, \frac{\partial f}{\partial y} = 0
        \)点中找到\(f(x, y)\)的最大值点或最小值点。
    \end{enumerate}
\end{conclusion}

\begin{conclusion}{求二元函数与三元函数的条件极值问题}{condition extremum problem}
    \begin{enumerate}
        \item 求函数\(z = f(x, y)\)在条件\(\varphi(x, y) = 0\)下的最大值或最小值
        \begin{enumerate}
            \item 化为无条件极值(最值),若从条件\(\varphi(x, y) = 0\)处可解出\(y = y(x)\)或\(x = x(y)\),则可将\(z = f(x, y)\)化为一元函数
            的最值问题
            \item 转化为拉格朗日函数，首先构造辅助函数(称为拉格朗日函数)\(
                F(x,y,\lambda)=f(x,y)+\lambda\varphi(x,y)
            \)然后求方程组\[\begin{cases}
                \frac{\partial F}{\partial x} = \frac{\partial f}{\partial x} + \lambda \frac{\partial \varphi}{\partial x} = 0 \\
                \frac{\partial F}{\partial y} = \frac{\partial f}{\partial y} + \lambda \frac{\partial \varphi}{\partial y} = 0 \\
                \frac{\partial F}{\partial \lambda} = \varphi(x, y) = 0
            \end{cases}\]的解，解\((x, y, \lambda)\)中的\((x, y)\)是在条件\(\varphi(x, y) = 0\)下的可能的极值点，最后比较这些可能极值点(若曲线\(\varphi(x, y)\)含端点，则还需再靠考察端点)的函数值，得到最大值点和最小值点。
       \end{enumerate}
        \item 求三元函数\(u = f(x, y, z)\)在条件\(\varphi(x, y, z) = 0, \psi(x, y, z) = 0\)下的最大值或最小值 \\
        构造拉格朗日函数\(F(x, y, z, \lambda, \mu) = f(x, y, z) + \lambda \varphi(x, y, z) + \mu \psi(x, y, z)\),然后求方程组\[\begin{cases}
            \frac{\partial F}{\partial x} = \frac{\partial f}{\partial x} + \lambda \frac{\partial \varphi}{\partial x} + \mu \frac{\partial \psi}{\partial x} = 0 \\
            \frac{\partial F}{\partial y} = \frac{\partial f}{\partial y} + \lambda \frac{\partial \varphi}{\partial y} + \mu \frac{\partial \psi}{\partial y} = 0 \\
            \frac{\partial F}{\partial z} = \frac{\partial f}{\partial z} + \lambda \frac{\partial \varphi}{\partial z} + \mu \frac{\partial \psi}{\partial z} = 0 \\
            \frac{\partial F}{\partial \lambda} = \varphi(x, y, z) = 0 \\
            \frac{\partial F}{\partial \mu} = \psi(x, y, z) = 0
            \end{cases}\]的解，解中的\((x, y, z)\)是在条件\(\varphi(x, y, z) = 0, \psi(x, y, z) = 0\)下的可能的极值点，最后比较这些可能
            极值点(若曲线\(\phi(x, y, z)\)或\(\psi(x, y, z)\)含端点，则还需再考察端点)的函数值，得到最大值点和最小值点。
    \end{enumerate}
\end{conclusion}
\subsection{梯度与方向导数}
\begin{mydef}{方向导数}{directional derivative}
    平面上过点\(M_0(x_0,y_0)\)以单位向量\(\mathbf l=(\cos\alpha,\cos\beta)\)为方向（其中\(\cos^2\alpha+\cos^2\beta=1\)）的直线参数方程为
    \[
        x = x_0 + t \cos \alpha, \quad y = y_0 + t \cos \beta
    \]
    \(z =f(x, y)\)限制在直线\(L\)上时,就变化为一元函数
    \[
        \varphi(t) = f(x_0 + t \cos \alpha, y_0 + t \cos \beta)
    \]
    若存在极限
    \[
        \varphi'(0) = \lim_{t \to 0} \frac{
            f(x_0 + t \cos \alpha, y_0 + t \cos \beta) - f(x_0, y_0)
        }{t}
    \]
    则称此极限为函数\(z = f(x, y)\)在点\(M_0(x_0, y_0)\)沿方向向量\(l\)的方向导数，记作\[
        \frac{\partial f(x_0, y_0)}{\partial l}  = 
        \left. \frac{\partial f}{\partial l} \right|_{(x_0, y_0)}
    \]
\end{mydef}
\begin{conclusion}{方向导数的存在性与计算公式}{directional derivative existence and formula}
    设\(f(x, y)\)在\(M(x_0, y_0)\)处可微，则\(f(x, y)\)在\(M_0\)点中找到任意方向向量\(l = (\cos \alpha, \cos \beta)\)的方向导数存在，且计算公式为\[
        \left. \frac{\partial f}{\partial l} \right|_{(x_0, y_0)} = 
        \frac{\partial f(x_0, y_0)}{\partial x} \cos \alpha + 
        \frac{\partial f(x_0, y_0)}{\partial y} \cos \beta
    \]
    若\(\alpha\)和\(\beta\)用极角表示，则有\[
        \left. \frac{\partial f(x_0, y_0)}{\partial l} \right|_{(x_0, y_0)} = 
        \frac{\partial f(x_0, y_0)}{\partial x} \cos \theta        + \frac{\partial f}{\partial y} \sin \theta
    \]
    对于三元函数的情形，设三元函数\(u = f(x, y, z)\)在点\(M_0(x_0, y_0, z_0)\)处可微，则三元函数\(u = f(x, y, z)\)在点\(M_0\)沿方向向量\(
    l = (\cos \alpha, \cos \beta, \cos \gamma)\)的方向导数存在，且计算公式为\[
        \left. \frac{\partial f}{\partial l} \right|_{(x_0, y_0, z_0)} = 
        \frac{\partial f(x_0, y_0, z_0)}{\partial x} \cos \alpha + 
        \frac{\partial f(x_0, y_0, z_0)}{\partial y} \cos \beta + 
        \frac{\partial f(x_0, y_0, z_0)}{\partial z} \cos \gamma
        \]
\end{conclusion}
\begin{conclusion}{梯度(向量)与方向导数的最大值}{gradient and directional derivative maximum}
    函数\(z = f(x, y)\)在点\(M_0\)的方向导数计算公式可改写为\[
    \begin{aligned}
        \left. \frac{\partial f(x_0, y_0)}{\partial l} \right|_{(x_0, y_0)} = 
        \left(
            \frac{\partial f(x_0, y_0)}{\partial x}, \frac{\partial f(x_0, y_0)}{\partial y}
        \right)
        \cdot (\cos \alpha, \cos \beta) \\
        = \mathbf{grad} f(x_0, y_0) \cdot \mathbf{l}
        = \left| \mathbf{grad} f(x_0, y_0) \right| \cos <\mathbf{grad} f(x_0, y_0), \mathbf{l}>
    \end{aligned}
        \]
    这里向量\(\mathbf{grad} f(x_0, y_0)\)是函数\(f\)在点\(M_0\)的梯度向量
    \(\frac{\partial f(x_0, y_0)}{\partial l}\)随\(\mathbf{l}\)的变化而变化，但当
    \(\mathbf{l} = \dfrac{\mathbf{grad} f(x_0, y_0)}{\left|
        \mathbf{grad} f(x_0, y_0)
    \right|} \)（梯度非零），即沿梯度方向时，方向导数达到最大值
    \(\left|
        \mathbf{grad} f(x_0, y_0)
    \right|\)
    因此梯度的方向是函数增长最快的方向，其模是最大方向导数；若梯度为零，则所有方向导数均为零，不存在唯一的最速增长方向。
\end{conclusion}
\subsection{多元函数的几何应用}
\begin{conclusion}{用隐式方程表示的曲面}{implicit equation surface}
    若空间曲面\(S\)的方程为\(F(x, y, z) = 0\)，\(M_0(x_0, y_0, z_0)\)是\(S\)上的一点，则\(S\)在\(M_0\)点处的切平面方程为
    \[
        \frac{\partial F(M_0)}{\partial x} (x - x_0) + 
        \frac{\partial F(M_0)}{\partial y} (y - y_0) + 
        \frac{\partial F(M_0)}{\partial z} (z - z_0) = 0
    \]
    同理，曲面\(S\)在点\(M_0(x_0, y_0, z_0)\)处的法线向量\(\mathbf{grad} F(M_0)\)
    是
    \[
        \left(
            \frac{\partial F(M_0)}{\partial x}, 
            \frac{\partial F(M_0)}{\partial y}, 
            \frac{\partial F(M_0)}{\partial z}
            \right)
    \]
    法线方程为
    \[
        \frac{x - x_0}{\frac{\partial F(M_0)}{\partial x}} = 
        \frac{y - y_0}{\frac{\partial F(M_0)}{\partial y}} = 
        \frac{z - z_0}{\frac{\partial F(M_0)}{\partial z}}
    \] 
    其中\(F(x, y ,z)\)在\((x_0, y_0, z_0)\)处有连续偏导数且
    \[
        \left[\frac{\partial F(M_0)}{\partial x}\right]^{2} + 
        \left[\frac{\partial F(M_0)}{\partial y}\right]^{2} + 
        \left[\frac{\partial F(M_0)}{\partial z}\right]^{2} \neq 0
    \]
\end{conclusion}
\begin{conclusion}{用显式方程表示的曲面}{explicit equation surface}
    若空间曲面\(S\)的方程为\(z=f(x,y)\)，且\(M_0(x_0,y_0,z_0)\)在曲面上，即\(z_0=f(x_0,y_0)\)，则曲面在点\(M_0\)处的切平面方程为
    \[
        z - z_0  = z - f(x_0, y_0)= \frac{\partial f}{\partial x}(x_0, y_0)(x - x_0) + 
        \frac{\partial f}{\partial y}(x_0, y_0)(y - y_0)
    \]
    法线方程为
    \[
        \frac{x - x_0}{\frac{\partial f(M_0)}{\partial x}} = 
        \frac{y - y_0}{\frac{\partial f(M_0)}{\partial y}} = 
        \frac{z - z_0}{-1}
    \]
    其中\(z = f(x, y)\)在\((x_0, y_0)\)处有连续偏导数.这里\(S\)在\(M_0\)的法向量为\(
        \pm (-f'_{x}(x_0, y_0), -f'_{y}(x_0, y_0), 1)
    \),
    单位法向量为\[
        \mathbf{n} = \dfrac{
            \pm(-f'_{x}(x_0, y_0), -f'_{y}(x_0, y_0), 1)
        }{
            \sqrt{
                \left[f'_{x}(x_0, y_0)\right]^{2} + 
                \left[f'_{y}(x_0, y_0)\right]^{2} + 1
            }
        } = (
            \cos \alpha, \cos \beta, \cos \gamma
        )\]
        这里\(\alpha, \beta, \gamma\)是所选单位法向量\(\mathbf n\)与\(x,y,z\)轴正向的方向角；它们不是切平面与坐标轴的夹角。\(\mathbf{n}\)的表达式中正负号的选取取决于\(\gamma\)是锐角还是钝角。
        即法向量是朝上还是朝下。朝上取正，朝下取负。
\end{conclusion}
\begin{conclusion}{用参数方程表示的空间曲线}{parameter equation space curve}
    若空间曲线\(\Gamma\)的参数方程为\(
        x = x(t), y = y(t), z = z(t) \quad (\alpha \leq t \leq \beta)
    \)又\(
        M_0(x_0, y_0, z_0) = (x(t_0), y(t_0), z(t_0)) \in \Gamma
    \)是\(\Gamma\)上的一点，则\(\Gamma\)在点\(M_0\)的切线方程为
    \[
        \frac{x - x_0}{x'(t_0)} = \frac{y - y_0}{y'(t_0)} = \frac{z - z_0}{z'(t_0)}
    \]
    法平面方程为
    \[
        x'(t_0)(x-x_0)+y'(t_0)(y-y_0)+z'(t_0)(z-z_0)=0.
    \]
    其中\(x(t), y(t), z(t)\)在\(t_0\)处有连续导数。
    且
    \(
        \left[x'(t_0)\right]^{2} + \left[y'(t_0)\right]^{2} + 
        \left[z'(t_0)\right]^{2} \neq 0
    \)
    其中
    \(
        x'(t_0), y'(t_0), z'(t_0)
    \)是曲线\(\Gamma\)在点\(M_0\)处的切向量.
\end{conclusion}
\begin{conclusion}{作为两曲面交线的空间曲线}{intersection of two curves}
    若空间曲线\(\Gamma\)的一般方程为\(
        \begin{cases}
            F(x, y, z) = 0 \\
            G(x, y, z) = 0
        \end{cases}
    \)是\(\Gamma\)上的一点，则\(\Gamma\)在点\(M_0\)处的切线方程为
    \[
        \begin{cases}
            \frac{\partial F(M_0)}{\partial x}(x - x_0) + 
            \frac{\partial F(M_0)}{\partial y}(y - y_0) + 
            \frac{\partial F(M_0)}{\partial z}(z - z_0) = 0 \\
            \frac{\partial G(M_0)}{\partial y}(y - y_0) + 
            \frac{\partial G(M_0)}{\partial x}(x - x_0) + 
            \frac{\partial G(M_0)}{\partial z}(z - z_0) = 0
        \end{cases}
    \]
    或
    \[
        \frac{x - x_0}{
            \left.
                \frac{\partial(F, G)}
                {\partial(y, z)}
            \right|_{M_0}
        }
        = \frac{y - y_0}{
            \left.
                \frac{\partial(F, G)}
                {\partial(z, x)}
            \right|_{M_0}
        }
        = \frac{z - z_0}{
            \left.
                \frac{\partial(F, G)}
                {\partial(x, y)}
            \right|_{M_0}
        }
    \]
    这里\(
        \frac{\partial(F, G)}
                {\partial(u, v)}
    \)表示
    \(
        \begin{vmatrix}
            \frac{\partial F}{\partial u} & \frac{\partial F}{\partial v} \\
            \frac{\partial G}{\partial u} & \frac{\partial G}{\partial v} 
        \end{vmatrix}
    \)
    法平面方程为
    \[
        \left.\frac{\partial(F, G)}{\partial(y, z)}\right|_{M_0} \cdot (x - x_0) + 
        \left.\frac{\partial(F, G)}{\partial(z, x)}\right|_{M_0} \cdot (y - y_0) + 
        \left.\frac{\partial(F, G)}{\partial(x, y)}\right|_{M_0} \cdot (z - z_0) = 0
    \]
    其中\(F(x, y ,z)\)和\(G(x, y ,z)\)在\((x_0, y_0, z_0)\)处有连续偏导数且
    \(
        \mathbf{grad}F(M_0) \times \mathbf{grad}G(M_0) \neq \mathbf{0}
    \)
    这里\(\Gamma\)在\(M_0\)的切向量是
    \[
        \begin{aligned}
             \mathbf{grad}F(M_0) \times \mathbf{grad}G(M_0)
        = \begin{vmatrix}
            \mathbf{i} & \mathbf{j} & \mathbf{k} \\

              \frac{\partial F}{\partial x} 
            & \frac{\partial F}{\partial y} 
            & \frac{\partial F}{\partial z} \\
            \frac{\partial G}{\partial x} 
            & \frac{\partial G}{\partial y} 
            & \frac{\partial G}{\partial z} \\

            \end{vmatrix} = \\
             \left(
                \left.\frac{\partial(F, G)}{\partial(y, z)}\right|_{M_0}, 
                \left.\frac{\partial(F, G)}{\partial(z, x)}\right|_{M_0}, 
                \left.\frac{\partial(F, G)}{\partial(x, y)}\right|_{M_0}
             \right)
        \end{aligned}
    \]
\end{conclusion}
\subsection{普适方法与技巧}
\begin{conclusion}{讨论二元函数$z = f(x,y)$可微的常用方法}{common method of differentiability}
    \begin{enumerate}
        \item 利用可微的定义，步骤为
        \begin{enumerate}
            \item 考察\(f'_{x}(x_0, y_0)\)和\(f'_{y}(x_0, y_0)\)是否存在，若不存在则\(f(x, y)\)在点\((x_0, y_0)\)处不可微。
            \item 若\(f'_{x}(x_0, y_0)\)和\(f'_{y}(x_0, y_0)\)存在，则考察
            \(
                \lim\limits_{\substack{
                    \varDelta x \to 0\\ \varDelta y \to 0}
                } \dfrac{
                    \varDelta z -
                    \left[f'_{x}(x_0, y_0)\varDelta x + f'_{y}(x_0, y_0)\varDelta y\right]
                }{\rho} = 0
            \)是否成立，其中\(\rho\)为\(\sqrt{\varDelta x^2 + \varDelta y^2}\)。\\
            若该极限等于\(0\)，则\(f(x, y)\)在点\((x_0, y_0)\)处可微；若极限不存在或不等于\(0\)，则不可微。
        \end{enumerate}
        \item 利用可微的必要条件证明\(f(x, y)\)是点\(x_0, y_0\)处不可微的，若能证明
        \(f(x, y)\)在点\((x_0, y_0)\)处不连续或不偏导,则\(f(x, y)\)在点\((x_0, y_0)\)处不可微。
        \item 利用可微的充分条件：证明\(f_x,f_y\)在该点的某邻域内存在，且在点\((x_0,y_0)\)处连续，即可推出\(f\)在该点可微。仅有函数连续不能推出可微。
    \end{enumerate}
\end{conclusion}

\begin{conclusion}{$k$次齐次函数}{homogeneous function}
    \(k\)次齐次函数是指形如
    \[
        f(ax_1, ax_2, \dots , ax_n) = a^k f(x_1, x_2, \dots , x_n)
    \]
    的函数，其中\(a\)为常数。
    \(k\)次齐次函数在点\((x_1, x_2, \dots , x_n)\)处的偏导数存在，则有
    \[
        \sum_{i = 1}^{n} x_i\frac{\partial f}{\partial x_i}(x_1, x_2, \dots , x_n)
        = kf(x_1, x_2, \dots , x_n)
    \]
\end{conclusion}


\subsection{基础过关题}

\textbf{1.} \textbf{（极限存在性判断）} 判断下列二重极限是否存在，若存在，求其值：
\[
\lim_{(x,y)\to(0,0)} \frac{xy^{2}}{x^{2}+y^{4}}.
\]

\boxed{\textbf{解}}
沿直线 $y=kx$：$\displaystyle\lim_{x\to 0}\frac{x(kx)^{2}}{x^{2}+(kx)^{4}} = \lim_{x\to 0}\frac{k^{2}x^{3}}{x^{2}+k^{4}x^{4}} = \lim_{x\to 0}\frac{k^{2}x}{1+k^{4}x^{2}} = 0$。

沿抛物线 $x=y^{2}$：$\displaystyle\lim_{y\to 0}\frac{y^{2}\cdot y^{2}}{(y^{2})^{2}+y^{4}} = \lim_{y\to 0}\frac{y^{4}}{2y^{4}} = \frac{1}{2}$。

沿不同路径得到不同极限值，故该二重极限\textbf{不存在}。

\vspace{6pt}

\textbf{2.} \textbf{（分段函数偏导数与可微性）} 设
\[
f(x,y)=\begin{cases}
\dfrac{xy}{\sqrt{x^{2}+y^{2}}}, & (x,y)\neq(0,0),\\[10pt]
0, & (x,y)=(0,0).
\end{cases}
\]
(1) 求 $f_{x}(0,0)$ 和 $f_{y}(0,0)$；\quad
(2) 讨论 $f(x,y)$ 在 $(0,0)$ 处的连续性与可微性。

\boxed{\textbf{解}}
(1) 由偏导数定义：
\[
f_{x}(0,0)=\lim_{\varDelta x\to 0}\frac{f(\varDelta x,0)-f(0,0)}{\varDelta x}
=\lim_{\varDelta x\to 0}\frac{0-0}{\varDelta x}=0.
\]
同理 $f_{y}(0,0)=0$。

(2) \textbf{连续性：} 由不等式 $|xy|\leqslant\frac12(x^{2}+y^{2})$，得
\[
|f(x,y)-0|=\frac{|xy|}{\sqrt{x^{2}+y^{2}}}\leqslant
\frac{\frac12(x^{2}+y^{2})}{\sqrt{x^{2}+y^{2}}}
=\frac12\sqrt{x^{2}+y^{2}}\to 0\quad((x,y)\to(0,0)).
\]
故 $\displaystyle\lim_{(x,y)\to(0,0)}f(x,y)=0=f(0,0)$，$f$ 在 $(0,0)$ 处\textbf{连续}。

\textbf{可微性：} 考察
\[
\frac{\varDelta z-[f_{x}(0,0)\varDelta x+f_{y}(0,0)\varDelta y]}{\rho}
=\frac{\frac{\varDelta x\,\varDelta y}{\sqrt{\varDelta x^{2}+\varDelta y^{2}}}}{\sqrt{\varDelta x^{2}+\varDelta y^{2}}}
=\frac{\varDelta x\,\varDelta y}{\varDelta x^{2}+\varDelta y^{2}},
\]
其中 $\rho=\sqrt{\varDelta x^{2}+\varDelta y^{2}}$。取 $\varDelta y=\varDelta x$，上式 $=\dfrac12\not\to 0$（$\rho\to 0$）。故 $f$ 在 $(0,0)$ 处\textbf{不可微}。

\boxed{\textbf{注}} 此例说明：偏导数存在且函数连续，仍不能保证可微。二元函数的可微性比可偏导性强得多。

\vspace{6pt}

\textbf{3.} \textbf{（全微分与近似计算）} 设 $z=\arctan\dfrac{y}{x}\;(x>0)$。
(1) 求 $\mathrm{d}z\big|_{(1,1)}$；
(2) 利用全微分近似计算 $\arctan\dfrac{1.02}{0.97}$ 的值（精确到 $10^{-3}$）。

\boxed{\textbf{解}}
(1)
\[
\frac{\partial z}{\partial x}=\frac{1}{1+(y/x)^{2}}\cdot\left(-\frac{y}{x^{2}}\right)=-\frac{y}{x^{2}+y^{2}},
\qquad
\frac{\partial z}{\partial y}=\frac{1}{1+(y/x)^{2}}\cdot\frac{1}{x}=\frac{x}{x^{2}+y^{2}}.
\]
在 $(1,1)$ 处：$z_{x}(1,1)=-\dfrac12$，$z_{y}(1,1)=\dfrac12$，故
\[
\mathrm{d}z\big|_{(1,1)} = -\frac12\,\mathrm{d}x + \frac12\,\mathrm{d}y.
\]

(2) 要近似计算 $\arctan\dfrac{1.02}{0.97}$，应取 $x_0=1$、$y_0=1$，并令 $x=0.97$、$y=1.02$，故 $\varDelta x=-0.03$、$\varDelta y=0.02$。于是
\[
z(1,1)=\arctan 1 = \frac{\pi}{4}\approx 0.7854.
\]
由全微分近似 $\varDelta z\approx\mathrm{d}z$：
\[
\varDelta z \approx -\frac12\times(-0.03) + \frac12\times0.02 = 0.015+0.01 = 0.025.
\]
故
\[
\arctan\frac{1.02}{0.97} \approx \frac{\pi}{4} + 0.025 \approx 0.7854 + 0.025 = 0.8104 \approx 0.810.
\]

\vspace{10pt}

\subsection{综合提高题}

\textbf{4.} \textbf{（链式法则·二阶偏导数）} 设 $z=f(u,v)$ 具有二阶连续偏导数，其中 $u=x^{2}+y^{2}$，$v=xy$。试用 $f$ 关于 $u,v$ 的偏导数表示 $\displaystyle\frac{\partial^{2}z}{\partial x\partial y}$。

\boxed{\textbf{解}}
一阶偏导数：
\[
\frac{\partial z}{\partial x}=f_{1}'\cdot\frac{\partial u}{\partial x}+f_{2}'\cdot\frac{\partial v}{\partial x}
=2x f_{1}' + y f_{2}'.
\]
对 $y$ 再求偏导：
\begin{align*}
\frac{\partial^{2}z}{\partial x\partial y}
&= \frac{\partial}{\partial y}\bigl(2x f_{1}' + y f_{2}'\bigr) \\
&= 2x\cdot\frac{\partial f_{1}'}{\partial y} + f_{2}' + y\cdot\frac{\partial f_{2}'}{\partial y}.
\end{align*}
由于 $f_{1}'$, $f_{2}'$ 仍通过 $u,v$ 依赖于 $x,y$，再次使用链式法则：
\begin{align*}
\frac{\partial f_{1}'}{\partial y} &= f_{11}''\cdot\frac{\partial u}{\partial y} + f_{12}''\cdot\frac{\partial v}{\partial y}
= 2y f_{11}'' + x f_{12}'', \\
\frac{\partial f_{2}'}{\partial y} &= f_{21}''\cdot\frac{\partial u}{\partial y} + f_{22}''\cdot\frac{\partial v}{\partial y}
= 2y f_{21}'' + x f_{22}''.
\end{align*}
由混合偏导数连续性知 $f_{12}'' = f_{21}''$，代入得
\begin{align*}
\frac{\partial^{2}z}{\partial x\partial y}
&= 2x(2y f_{11}'' + x f_{12}'') + f_{2}' + y(2y f_{21}'' + x f_{22}'') \\
&= 4xy f_{11}'' + 2(x^{2}+y^{2}) f_{12}'' + xy f_{22}'' + f_{2}'.
\end{align*}

\boxed{\textbf{注}} 对于含抽象函数的高阶链式求导，关键是明确每次求导作用于哪个中间变量；画出变量依赖关系可避免遗漏交叉项。$f_{1}'$ 和 $f_{2}'$ 仍以 $(u,v)$ 为变量，求导时必须再次展开。

\vspace{6pt}

\textbf{5.} \textbf{（隐函数求导·二阶）} 设方程 $x^{2}+y^{2}+z^{2}-3xyz=0$ 确定了隐函数 $z=z(x,y)$，且 $z(1,1)=1$。求 $\displaystyle\frac{\partial z}{\partial x}\bigg|_{(1,1)}$ 和 $\displaystyle\frac{\partial^{2}z}{\partial x^{2}}\bigg|_{(1,1)}$。

\boxed{\textbf{解}}
\textbf{法一（隐函数求导公式）：} 令 $F(x,y,z)=x^{2}+y^{2}+z^{2}-3xyz$，则
\[
\frac{\partial z}{\partial x} = -\frac{F_{x}}{F_{z}}
= -\frac{2x-3yz}{2z-3xy}.
\]
在 $(1,1,1)$ 处：$\displaystyle\frac{\partial z}{\partial x}\bigg|_{(1,1)}
= -\frac{2-3}{2-3} = -1$。

\textbf{法二（方程两边直接求导，更推荐）：} 将 $x^{2}+y^{2}+z^{2}-3xyz=0$ 两边对 $x$ 求偏导（$y$ 视为常数）：
\[
2x + 2z z_{x} - 3y(z + x z_{x}) = 0,
\]
即 $(2z-3xy)z_{x} + 2x - 3yz = 0$。代入 $(1,1,1)$：$(2-3)z_{x}+2-3=0$，得 $z_{x}=-1$。

对 $x$ 再次求偏导：
\[
2 + 2(z_{x})^{2} + 2z z_{xx} - 3y(z_{x}+z_{x}+x z_{xx}) = 0,
\]
即 $2 + 2(z_{x})^{2} + 2z z_{xx} - 6y z_{x} - 3xy z_{xx} = 0$。
代入 $(1,1,1)$ 及 $z_{x}=-1$：
\[
2 + 2(1) + 2\cdot 1\cdot z_{xx} - 6\cdot 1\cdot(-1) - 3\cdot 1\cdot 1\cdot z_{xx} = 0,
\]
\[
2 + 2 + 2z_{xx} + 6 - 3z_{xx} = 0 \;\Longrightarrow\; 10 - z_{xx} = 0,
\]
故 $\displaystyle\frac{\partial^{2}z}{\partial x^{2}}\bigg|_{(1,1)} = 10$。

\boxed{\textbf{注}} 法二直接对恒等式 $F(x,y,z(x,y))=0$ 求导，在含参讨论和二阶导计算中通常更便于核对各项。

\vspace{6pt}

\textbf{6.} \textbf{（无条件极值）} 求函数 $f(x,y)=x^{3}+y^{3}-3x^{2}-3y^{2}-9x$ 的所有极值点，并判定极值类型。

\boxed{\textbf{解}}
求驻点：
\[
\begin{cases}
f_{x}=3x^{2}-6x-9=3(x^{2}-2x-3)=3(x-3)(x+1)=0,\\[4pt]
f_{y}=3y^{2}-6y=3y(y-2)=0.
\end{cases}
\]
解得 $x=-1$ 或 $x=3$；$y=0$ 或 $y=2$。故四个驻点：
\[
P_{1}(-1,0),\quad P_{2}(-1,2),\quad P_{3}(3,0),\quad P_{4}(3,2).
\]

二阶偏导：$f_{xx}=6x-6$，$f_{yy}=6y-6$，$f_{xy}=0$。
判别式 $\varDelta = f_{xx}f_{yy}-(f_{xy})^{2} = (6x-6)(6y-6)$。

\begin{itemize}
    \item $P_{1}(-1,0)$：$f_{xx}=-12<0$，$\varDelta=(-12)(-6)=72>0$ $\Rightarrow$ \textbf{极大值}，
    $f(-1,0)=-1+0-3-0+9=5$。
    \item $P_{2}(-1,2)$：$\varDelta=(-12)(6)=-72<0$ $\Rightarrow$ \textbf{非极值点}（鞍点）。
    \item $P_{3}(3,0)$：$\varDelta=(12)(-6)=-72<0$ $\Rightarrow$ \textbf{非极值点}（鞍点）。
    \item $P_{4}(3,2)$：$f_{xx}=12>0$，$\varDelta=(12)(6)=72>0$ $\Rightarrow$ \textbf{极小值}，
    $f(3,2)=27+8-27-12-27=-31$。
\end{itemize}

\boxed{\textbf{注}} 此题的巧妙之处在于 $f_{xy}=0$ 使判别式简化为两个单变量因子的乘积，每个驻点的类型由 $x$ 和 $y$ 分量的曲率方向独立决定——这是可分离变量的极值问题的典型特征。

\vspace{6pt}

\textbf{7.} \textbf{（拉格朗日乘数法）} 求函数 $f(x,y)=x^{2}+y^{2}$ 在椭圆 $x^{2}+xy+y^{2}=3$ 上的最大值和最小值。

\boxed{\textbf{解}}
约束曲线为有界闭集，$f$ 连续，故最值存在。构造拉格朗日函数：
\[
L(x,y,\lambda)=x^{2}+y^{2}+\lambda(x^{2}+xy+y^{2}-3).
\]
令
\[
\begin{cases}
L_{x}=2x+\lambda(2x+y)=0,\\[4pt]
L_{y}=2y+\lambda(x+2y)=0,\\[4pt]
L_{\lambda}=x^{2}+xy+y^{2}-3=0.
\end{cases}
\]
由前两式得线性方程组：
\[
\begin{pmatrix}
2(1+\lambda) & \lambda \\[2pt]
\lambda & 2(1+\lambda)
\end{pmatrix}
\begin{pmatrix} x \\ y \end{pmatrix}
= \begin{pmatrix} 0 \\ 0 \end{pmatrix}.
\]
存在非零解 $(x,y)\neq(0,0)$（否则 $x=y=0$ 不满足约束），故系数矩阵行列式为零：
\[
4(1+\lambda)^{2}-\lambda^{2}=0 \;\Longrightarrow\; (2+2\lambda-\lambda)(2+2\lambda+\lambda)=0,
\]
即 $(\lambda+2)(3\lambda+2)=0$，得 $\lambda=-2$ 或 $\lambda=-\dfrac{2}{3}$。

\begin{itemize}
    \item 当 $\lambda=-2$ 时，第一式给出 $2x-2(2x+y)=0\Rightarrow x+y=0\Rightarrow y=-x$。
    代入约束：$x^{2}-x^{2}+x^{2}=x^{2}=3\Rightarrow x=\pm\sqrt{3}$。
    得 $(\sqrt{3},-\sqrt{3})$ 和 $(-\sqrt{3},\sqrt{3})$。$f=3+3=6$。

    \item 当 $\lambda=-\dfrac{2}{3}$ 时，第一式给出 $2x-\dfrac{2}{3}(2x+y)=0\Rightarrow \dfrac{2}{3}x-\dfrac{2}{3}y=0\Rightarrow x=y$。
    代入约束：$3x^{2}=3\Rightarrow x=\pm1$。
    得 $(1,1)$ 和 $(-1,-1)$。$f=1+1=2$。
\end{itemize}
比较四个候选点的函数值：最大值为 $6$，最小值为 $2$。

\boxed{\textbf{注}} 此题等价于求椭圆 $x^{2}+xy+y^{2}=3$ 上点到原点距离平方的极值。注意到约束可配方为 $\dfrac{1}{2}[(x+y)^{2}+3(x-y)^{2}]=3$，正定二次型保证了约束为有界闭曲线，连续函数在此上必有最值。

\vspace{10pt}

\subsection{真题风格与综合训练}

\textbf{8.}（\boxed{\textbf{真题风格}}）\textbf{（链式变换与偏微分方程）} 设 $u=f(x,y)$ 具有二阶连续偏导数。作自变量变换 $\xi=x+ay$，$\eta=x-ay$（$a\neq0$ 为常数）。试用 $u$ 关于 $\xi,\eta$ 的二阶偏导数表示 $\displaystyle\frac{\partial^{2}u}{\partial x^{2}}-\frac{1}{a^{2}}\frac{\partial^{2}u}{\partial y^{2}}$。

\boxed{\textbf{解}}
由链式法则：
\[
\frac{\partial u}{\partial x}
= \frac{\partial u}{\partial\xi}\cdot\frac{\partial\xi}{\partial x}
  +\frac{\partial u}{\partial\eta}\cdot\frac{\partial\eta}{\partial x}
= u_{\xi}+u_{\eta}.
\]
\[
\frac{\partial u}{\partial y}
= \frac{\partial u}{\partial\xi}\cdot\frac{\partial\xi}{\partial y}
  +\frac{\partial u}{\partial\eta}\cdot\frac{\partial\eta}{\partial y}
= a u_{\xi} - a u_{\eta} = a(u_{\xi}-u_{\eta}).
\]

二阶偏导（注意 $u_{\xi},u_{\eta}$ 仍以 $(\xi,\eta)$ 为变量）：
\begin{align*}
\frac{\partial^{2}u}{\partial x^{2}}
&= \frac{\partial}{\partial x}(u_{\xi}+u_{\eta})
 = (u_{\xi\xi}+u_{\xi\eta}) + (u_{\eta\xi}+u_{\eta\eta}) \\
&= u_{\xi\xi} + 2u_{\xi\eta} + u_{\eta\eta}.
\end{align*}
\begin{align*}
\frac{\partial^{2}u}{\partial y^{2}}
&= \frac{\partial}{\partial y}\bigl[a(u_{\xi}-u_{\eta})\bigr]
 = a\bigl[(u_{\xi\xi}\cdot a + u_{\xi\eta}\cdot(-a)) - (u_{\eta\xi}\cdot a + u_{\eta\eta}\cdot(-a))\bigr] \\
&= a\bigl[a u_{\xi\xi} - a u_{\xi\eta} - a u_{\eta\xi} + a u_{\eta\eta}\bigr] \\
&= a^{2}(u_{\xi\xi} - 2u_{\xi\eta} + u_{\eta\eta}).
\end{align*}

代入目标表达式：
\begin{align*}
\frac{\partial^{2}u}{\partial x^{2}}-\frac{1}{a^{2}}\frac{\partial^{2}u}{\partial y^{2}}
&= (u_{\xi\xi}+2u_{\xi\eta}+u_{\eta\eta}) - \frac{1}{a^{2}}\cdot a^{2}(u_{\xi\xi}-2u_{\xi\eta}+u_{\eta\eta}) \\
&= u_{\xi\xi}+2u_{\xi\eta}+u_{\eta\eta} - u_{\xi\xi}+2u_{\xi\eta}-u_{\eta\eta} \\
&= 4u_{\xi\eta}.
\end{align*}

\boxed{\textbf{注}} 这是波动方程 $u_{xx}-a^{-2}u_{yy}=0$（将 $y$ 视为时间变量）的标准化简。变换后方程为 $u_{\xi\eta}=0$，其通解为 $u=\varphi(\xi)+\psi(\eta)=\varphi(x+ay)+\psi(x-ay)$，即经典的 d'Alembert 公式。此类变量替换是考研高频考点。

\vspace{6pt}

\textbf{9.}（\boxed{\textbf{真题风格}}）\textbf{（隐函数组与空间曲线）} 设方程组
\[
\begin{cases}
x^{2}+y^{2}+z^{2}=2,\\
x+y+z=0
\end{cases}
\]
在点 $(1,-1,0)$ 附近确定了隐函数 $y=y(x)$, $z=z(x)$。求
(1) $\displaystyle\frac{\mathrm{d}y}{\mathrm{d}x}$, $\displaystyle\frac{\mathrm{d}z}{\mathrm{d}x}$ 在 $x=1$ 处的值；
(2) 曲线在点 $(1,-1,0)$ 处的切线方程。

\boxed{\textbf{解}}
(1) 对方程组两边对 $x$ 求导（$y,z$ 是 $x$ 的函数）：
\[
\begin{cases}
2x + 2y y' + 2z z' = 0,\\
1 + y' + z' = 0.
\end{cases}
\]
代入 $(x,y,z)=(1,-1,0)$：
\[
\begin{cases}
2 + 2(-1)y' + 2\cdot 0\cdot z' = 0,\\
1 + y' + z' = 0.
\end{cases}
\;\Longrightarrow\;
\begin{cases}
2 - 2y' = 0,\\
1 + y' + z' = 0.
\end{cases}
\]
由第一式得 $y'=1$，代入第二式得 $1+1+z'=0\Rightarrow z'=-2$。
故 $\displaystyle\frac{\mathrm{d}y}{\mathrm{d}x}\Big|_{x=1}=1$，$\displaystyle\frac{\mathrm{d}z}{\mathrm{d}x}\Big|_{x=1}=-2$。

(2) 空间曲线的切向量为 $(1, y', z')|_{(1,-1,0)} = (1, 1, -2)$。
切线方程（参数式）：
\[
\frac{x-1}{1} = \frac{y+1}{1} = \frac{z-0}{-2}.
\]
即
\[
x-1 = y+1 = -\frac{z}{2}.
\]

\boxed{\textbf{注}} 本题的方程组 $F(x,y,z)=0$, $G(x,y,z)=0$ 确定了空间曲线（球面与平面的交线——一个圆）。对 $x$ 求导时，$y$ 和 $z$ 均视为 $x$ 的函数，得到关于 $y',z'$ 的线性方程组。这与雅可比行列式公式 $\dfrac{\mathrm{d}y}{\mathrm{d}x}=-\dfrac{\partial(F,G)/\partial(x,z)}{\partial(F,G)/\partial(y,z)}$ 等价，但直接求导在数值代入时更简便。

\vspace{6pt}

\textbf{10.}（\boxed{\textbf{真题风格}}）\textbf{（条件极值·几何应用）} 设 $a,b,c>0$，求内接于椭球面 $\dfrac{x^{2}}{a^{2}}+\dfrac{y^{2}}{b^{2}}+\dfrac{z^{2}}{c^{2}}=1$ 且各面平行于坐标面的长方体（第一卦限部分顶点坐标 $(x,y,z)$ 均为正）的最大体积。

\boxed{\textbf{解}}
由对称性，设长方体在第一卦限的顶点为 $(x,y,z)$（$x,y,z>0$），则长方体边长分别为 $2x$, $2y$, $2z$，体积 $V=8xyz$。
约束条件：$\dfrac{x^{2}}{a^{2}}+\dfrac{y^{2}}{b^{2}}+\dfrac{z^{2}}{c^{2}}=1$。

取对数简化计算：最大化 $V$ 等价于最大化 $\ln V=\ln 8+\ln x+\ln y+\ln z$。
构造拉格朗日函数：
\[
L(x,y,z,\lambda)=\ln x+\ln y+\ln z+\lambda\!\left(\frac{x^{2}}{a^{2}}+\frac{y^{2}}{b^{2}}+\frac{z^{2}}{c^{2}}-1\right).
\]
令
\[
\begin{cases}
L_{x}=\dfrac{1}{x}+2\lambda\dfrac{x}{a^{2}}=0,\\[8pt]
L_{y}=\dfrac{1}{y}+2\lambda\dfrac{y}{b^{2}}=0,\\[8pt]
L_{z}=\dfrac{1}{z}+2\lambda\dfrac{z}{c^{2}}=0,\\[8pt]
L_{\lambda}=\dfrac{x^{2}}{a^{2}}+\dfrac{y^{2}}{b^{2}}+\dfrac{z^{2}}{c^{2}}-1=0.
\end{cases}
\]
由前三式分别乘以 $x,y,z$，可得
\[
\frac{x^2}{a^2}=\frac{y^2}{b^2}=\frac{z^2}{c^2}=-\frac{1}{2\lambda}.
\]
代入约束：
\[
3\cdot\frac{x^{2}}{a^{2}}=1 \;\Longrightarrow\; \frac{x^{2}}{a^{2}}=\frac{1}{3},
\]
得 $x=\dfrac{a}{\sqrt{3}}$，$y=\dfrac{b}{\sqrt{3}}$，$z=\dfrac{c}{\sqrt{3}}$。

（由问题的实际意义，最大值存在且唯一驻点必为最大值点，无需验证二阶条件。）

最大体积：
\[
V_{\max}=8xyz=8\cdot\frac{a}{\sqrt{3}}\cdot\frac{b}{\sqrt{3}}\cdot\frac{c}{\sqrt{3}}=\frac{8abc}{3\sqrt{3}}.
\]

\boxed{\textbf{注}} 椭球内接最大长方体的体积为 $\dfrac{8abc}{3\sqrt{3}}$。对连乘积取对数后再使用 Lagrange 乘数法，可以把一阶条件迅速化为 $\dfrac{x^{2}}{a^{2}}=\dfrac{y^{2}}{b^{2}}=\dfrac{z^{2}}{c^{2}}$。
